
TODO: extra energy in battery is added to the profit as if it was sold at market price at the end of the simulation: $terminal_delta_value = (SoC_end - SoC_start) * eta_out * terminal_price$

\section{Decision and Economic Analysis}
This section applies a quantitative, model-based techno-economic analysis aligned with the research question: whether probabilistic forecasting improves economically relevant \gls{bess} dispatch decisions under physical and market constraints. The analytical framework combines (i) continuous linear optimization for ex-ante decision support, (ii) rolling-horizon backtesting with ex-post settlement, and (iii) comparative performance diagnostics across forecast quantiles and benchmark strategies. Hence, the methodological focus is not interpretive text analysis, but statistical-economic evaluation of realized cashflows, penalties, and risk-adjusted operational outcomes.

\subsection{Dispatch Optimization}
\label{subsec:dispatch-optimization}
The dispatch layer is formulated as a rolling continuous linear optimization problem that maximizes expected operating value subject to battery physics, inverter limits, and market participation constraints. In each optimization window, the decision vector includes \gls{da} charging/discharging setpoints, directional reserve commitments, auxiliary power variables, slack variables, and state trajectories.

Let $t\in\{1,\dots,T\}$ index hourly time steps. The soft-penalty objective is:
\begin{equation}\label{eq:milp-objective-soft}
    \max \; \hat{\Pi}^{\mathrm{LP}}
    =
    \sum_{t=1}^{T}\left(\hat{\pi}_{t}^{\mathrm{DA/ID}}+\hat{\pi}_{t}^{\mathrm{aFRR}}\right)
    - \lambda \sum_{t=1}^{T}\left(s_{t}^{+}+s_{t}^{-}\right),
\end{equation}
where $s_t^+\ge 0$ and $s_t^-\ge 0$ are slack variables for upward/downward \gls{soc} safety violations, and $\lambda>0$ is the severe imbalance-penalty coefficient. This term enforces risk awareness without rendering extreme forecast scenarios trivially infeasible.

The quantile-based reserve-feasibility constraints are implemented as soft constraints:
\begin{equation}\label{eq:soft-p90-pos}
    SoC_t - \sum_{b\in\mathcal{B}}\Big(\hat{p}^{+}_{\mathrm{acc},b,t}\,\hat{r}^{+}_{\tau,t}\,P^{+}_{b,t}\,\Delta t/\eta_{\text{out}}\Big) + s_t^{+}
    \ge SoC_{\min},
\end{equation}
\begin{equation}\label{eq:soft-p90-neg}
    SoC_t + \sum_{b\in\mathcal{B}}\Big(\hat{p}^{-}_{\mathrm{acc},b,t}\,\hat{r}^{-}_{\tau,t}\,P^{-}_{b,t}\,\Delta t\,\eta_{\text{in}}\Big) - s_t^{-}
    \le SoC_{\max}.
\end{equation}
Here, $\tau$ denotes the selected activation-rate forecast quantile used in the quantile sweep, for example $\tau \in \{0.50,\dots,0.99\}$.
Thus, potential quantile-based headroom/footroom violations are explicitly priced in the objective rather than hidden in infeasibility artifacts.

\subsection{Backtesting}
Backtesting is performed in a causal rolling-horizon setting. At decision time $t$, forecasts available at $t$ are used to solve the continuous dispatch problem over horizon $H$ (here: 48 hours), but only the first control action is implemented:
\begin{equation}\label{eq:rolling-policy}
    \mathbf{u}^{*}_{t:t+H-1}
    =
    \arg\max_{\mathbf{u}} \hat{\Pi}_{t:t+H-1}(\mathbf{u}\mid\mathcal{F}_t),
    \qquad
    \mathbf{u}^{\mathrm{exec}}_{t}=\mathbf{u}^{*}_{t}.
\end{equation}
After market clearing and physical settlement, the state is updated and the optimization is repeated at $t+1$. This reproduces operational receding-horizon control and avoids look-ahead bias.

Daily \gls{da} gate closure is represented by lock-in logic: previously committed \gls{da} schedules for already-closed delivery intervals are fixed in subsequent re-optimizations. Hence, only non-closed periods remain decision variables, while historical commitments are treated as binding constraints.

Realized ex-post settlement is evaluated on true market outcomes:
\begin{equation}\label{eq:realized-hourly-pnl}
    \pi^{\mathrm{real}}_{t}
    =
    \pi^{\mathrm{DA/ID}}_{t}
    +
    \pi^{\mathrm{aFRR}}_{t}
    -
    C_{pen,t}
    -
    C_{tx,t}
    -
    C^{\mathrm{deg}}_{t}.
\end{equation}
The cumulative realized strategy value is
\begin{equation}\label{eq:realized-total-pnl}
    \Pi^{\mathrm{real}}=\sum_{t=1}^{T}\pi^{\mathrm{real}}_t + V_{\mathrm{term}}.
\end{equation}

For benchmarking, two reference policies are used:
\begin{align}
    \Pi^{\mathrm{naive}}  & : \text{forecast-driven benchmark without advanced probabilistic structure},          \\
    \Pi^{\mathrm{oracle}} & : \text{perfect-foresight upper bound using realized prices/activations in planning}.
\end{align}
The oracle is interpreted as an informational upper bound, not as an implementable trading policy.


This section is organized as follows: physical feasibility and \gls{soc} accounting first, then \gls{da}/\gls{id} and \gls{afrr} profit models (ex-post and ex-ante), followed by the master objective and assumptions.
\begin{longtable}{@{}>{\raggedright\arraybackslash}p{0.34\textwidth} >{\raggedright\arraybackslash}p{0.46\textwidth} >{\raggedright\arraybackslash}p{0.14\textwidth}@{}}
    \caption{Compact nomenclature (symbols used in Chapter~\ref{ch:analysis}).}\label{tab:nomenclature}                                                                                                                                                                                                                                                                                                                                                                                  \\
    \toprule
    \textbf{Symbol}                                                                                                                                                                                                                                                                                           & \textbf{Description}                                                                                                               & \textbf{Unit}                       \\
    \midrule
    \endfirsthead
    \multicolumn{3}{l}{\textit{Table \ref{tab:nomenclature} continued}}                                                                                                                                                                                                                                                                                                                                                                                                                  \\
    \toprule
    \textbf{Symbol}                                                                                                                                                                                                                                                                                           & \textbf{Description}                                                                                                               & \textbf{Unit}                       \\
    \midrule
    \endhead
    \midrule
    \multicolumn{3}{r}{\textit{Continued on next page}}                                                                                                                                                                                                                                                                                                                                                                                                                                  \\
    \endfoot
    \bottomrule
    \endlastfoot
    \begin{tabular}[t]{@{}l@{}}$SoC_{0/t/T},\,SoC_{\min/\max}$\end{tabular}                                                                                                                                                                                                                                   & State of charge variable and lower/upper/initial/terminal bounds ($SoC_t$: end-of-interval after modeled trading/activation flows) & MWh                                 \\
    $E_{\text{cap}},\,P_{\text{max}},\,\Delta t,\,T$                                                                                                                                                                                                                                                          & Battery energy/power ratings and horizon indexing                                                                                  & MWh, MW, h, -                       \\
    \begin{tabular}[t]{@{}l@{}}$E_{\text{recv},t},\,E_{\text{sent},t}$\\$E_{\text{recv},t}^{\text{plan}},\,E_{\text{sent},t}^{\text{plan}}$\\$E^{DA/ID}_{\mathrm{recv},t},\,E^{DA/ID}_{\mathrm{sent},t}$\end{tabular}                                                                                         & DA/ID charging/discharging energies (settled ex-post and planned ex-ante)                                                          & MWh                                 \\
    $E_{\text{act},t},\,E_{\text{act},t}^{+},\,E_{\text{act},t}^{-}$                                                                                                                                                                                                                                          & Signed and split activation energies (upward/downward components)                                                                  & MWh                                 \\
    $E_{\text{recv},t}^{\text{SoC}},\,E_{\text{sent},t}^{\text{SoC}}$                                                                                                                                                                                                                                         & Total physical SoC-relevant charging/discharging energies                                                                          & MWh                                 \\
    $\eta_{\text{in}},\,\eta_{\text{out}}$                                                                                                                                                                                                                                                                    & Charging/discharging efficiencies                                                                                                  & -                                   \\
    $\lambda_{\text{mkt},t},\,\lambda_{\text{clr},t}$                                                                                                                                                                                                                                                         & Realized DA/ID market price and aFRR activation clearing price                                                                     & €/MWh                               \\
    $\hat{\lambda}_{\text{mkt},t},\,\hat{\lambda}_{\text{clr},t}$                                                                                                                                                                                                                                             & Forecasted DA/ID market price and forecasted aFRR clearing price                                                                   & €/MWh                               \\
    $C_{\text{deg}}$                                                                                                                                                                                                                                                                                          & Degradation cost parameter                                                                                                         & €/MWh                               \\
    $\Delta P_{j}^{d}$                                                                                                                                                                                                                                                                                        & Relative bid spread interval $j$ for direction $d$                                                                                 & €/MWh                               \\
    $J,\,j$                                                                                                                                                                                                                                                                                                   & Set of discrete bid spread intervals and index                                                                                     & -                                   \\
    \begin{tabular}[t]{@{}l@{}}$P_{\text{charge},t},\,P_{\text{discharge},t}$\\$R_t^{+},\,R_t^{-},\,R_{\max}$\\$P_{\text{bid},t},\,P_{\text{bid},t}^{\uparrow},\,P_{\text{bid},t}^{\downarrow}$\\$P_{\text{bid,min/max}},\,\Delta P_{\text{bid}}$\end{tabular}                                                                  & Continuous charging/discharging, reserve, and bid-volume variables                                                                & MW                                  \\
    $m_{j}^{d}$                                                                                                                                                                                                                                                                                               & Scalar midpoint of the spread interval $j$                                                                                         & €/MWh                               \\
    $\hat{p}_{\text{acc}, j, t}^{d}$                                                                                                                                                                                                                                                                          & Forecasted probability of bid acceptance for interval $j$                                                                          & \%                                  \\
    $\tau,\,\hat{r}_{\tau,t}^{+},\,\hat{r}_{\tau,t}^{-}$                                                                                                                                                                                                                                                      & Selected activation-rate quantile and directional activation-rate forecasts used in the quantile sweep                             & -, \%                               \\
    $s_t^{+},\,s_t^{-},\,\lambda$                                                                                                                                                                                                                                                                             & Soft-constraint slack variables and penalty coefficient                                                                            & MWh, €/MWh                          \\
    $\mathcal{B},\,b,\,P_{b,t}^{+},\,P_{b,t}^{-}$                                                                                                                                                                                                                                                             & Bid-spread interval set, interval index, and directional bid volumes in the soft reserve-feasibility constraints                   & -, -, MW                            \\
    $V_{\text{bid}, t}^{d}$                                                                                                                                                                                                                                                                                   & Submitted bid volume for direction $d$                                                                                             & MW                                  \\
    $C_{\text{bid},t},\,C_{\text{bid},t}^{\uparrow,\downarrow},\,\hat{C}_{\text{bid},t}$                                                                                                                                                                                                                      & Realized directional and forecasted aFRR capacity prices                                                                           & €/MW/h                              \\
    $C_{\text{pos},t},\,C_{\text{neg},t}$                                                                                                                                                                                                                                                                     & Non-energy activation costs for upward/downward activation                                                                         & €                                   \\
    $\pi_t^{\text{chg}},\,\pi_t^{\text{dis}},\,\pi_t^{\text{DA/ID}}$                                                                                                                                                                                                                                          & Realized DA/ID cashflow components and total                                                                                       & €                                   \\
    $\pi_t^{\text{cap}},\,\pi_t^{\text{act}},\,\pi_t^{\text{aFRR}}$                                                                                                                                                                                                                                           & Realized aFRR capacity, activation, and total profit components                                                                    & €                                   \\
    $\hat{\pi}_t^{\text{DA/ID}},\,\hat{\pi}_t^{\text{aFRR}}$                                                                                                                                                                                                                                                  & Expected DA/ID and aFRR profit terms                                                                                               & €                                   \\
    $\mathbf{X}_t,\,g_{\text{bid}}(\cdot),\,g_{\text{clr}}(\cdot),\,g_{\uparrow}(\cdot),\,g_{\downarrow}(\cdot)$                                                                                                                                                                                              & Feature vector and forecast mappings                                                                                               & -                                   \\
    $C_{\text{marg},t}^{\uparrow},\,C_{\text{marg},t}^{\downarrow}$                                                                                                                                                                                                                                           & Directional marginal activation thresholds                                                                                         & €/MWh                               \\
    $P_{\text{proc},t}^{\uparrow,\downarrow},\,E_{\text{act},t}^{\uparrow,\downarrow}$                                                                                                                                                                                                                        & Procured and activated system-level reserve quantities                                                                             & MW, MWh                             \\
    $\hat{\Pi}_{\text{tot}},\,\Pi^{\mathrm{real}}$                                                                                                                                                                                                                                                            & Total expected and realized objective values                                                                                       & €                                   \\
    $C_{tx,t}$, $c_{tx}$, $C_{pen,t}$                                                                                                                                                                                                                                                                         & Transaction cost, specific transaction fee, and non-delivery penalties at time $t$                                                 & \texteuro, \texteuro/MWh, \texteuro \\
    $C^{\mathrm{deg}}_t,\,C^{\mathrm{tx}},\,C^{\mathrm{pen}},\,C^{\mathrm{deg}}$                                                                                                                                                                                                                              & Period-specific degradation cost and aggregate transaction, penalty, and degradation costs over the backtest                       & \texteuro                           \\
    $CF_{net,t}$                                                                                                                                                                                                                                                                                              & Net liquid cashflow at time $t$ (excluding non-cash degradation)                                                                   & \texteuro                           \\
    $Cash_0$                                                                                                                                                                                                                                                                                                  & Initial cash balance                                                                                                               & \texteuro                           \\
    $V_{\mathrm{term}}$                                                                                                                                                                                                                                                                                       & Terminal Mark-to-Market value of the remaining sellable inventory                                                                  & \texteuro                           \\
    $\text{max\_cap\_req}$                                                                                                                                                                                                                                                                                    & Peak cumulative drawdown / maximum working capital required                                                                        & \texteuro                           \\
    $\text{EFC}$                                                                                                                                                                                                                                                                                              & Equivalent Full Cycles (internal chemical throughput)                                                                              & -                                   \\
    $ROI_{cap}$                                                                                                                                                                                                                                                                                               & Return on maximum capital requirement                                                                                              & \%                                  \\
\end{longtable}

The complete list of symbols is provided in Appendix~\ref{sec:app-nomenclature}.

\subsubsection{Model Scope and Timeline}
The optimization horizon is indexed by $t=1,\dots,T$ with fixed hourly granularity $\Delta t=1$ h.
The formulation separates ex-post (realized) settlement accounting from ex-ante expected-value decision support.
In the rolling 48-hour receding-horizon implementation, the optimization is solved again for each daily auction using the newest forecasts.
Only the decisions for the next auction are executed, while the remaining 48-hour plan is overwritten on the following day; therefore, energy left in the battery is carried forward as realized \gls{soc} and valued through the terminal mark-to-market term instead of being forced to be sold at the horizon end.

\subsubsection{Numerical Parameterization}
Table~\ref{tab:parameterization} summarizes the baseline numerical setup used in the current implementation.

\begin{table}[H]
    \centering
    \caption{Baseline model parameterization and open parameter definitions.}\label{tab:parameterization}
    \begin{tabularx}{\textwidth}{@{}>{\raggedright\arraybackslash}p{0.16\textwidth} >{\raggedright\arraybackslash}p{0.13\textwidth} >{\raggedright\arraybackslash}p{0.11\textwidth} >{\raggedright\arraybackslash}X@{}}
        \toprule
        \textbf{Symbol}             & \textbf{Value}             & \textbf{Unit} & \textbf{Rationale}                                                                                \\
        \midrule
        $E_{\text{cap}}$            & 20.0                       & MWh           & Matches commercial BESS projects, such as Mühlacker (10 MW / 20 MWh) and Marbach (15 MW / 22 MWh) \\
        $P_{\text{max}}$            & 10.0                       & MW            & Matches commercial BESS projects, such as Mühlacker (10 MW / 20 MWh)                              \\
        $\eta_{\text{rt}}$          & 0.90                       & -             & AC-to-AC round-trip efficiency                                                                    \\
        $\eta_{\text{in}}$          & 0.9487                     & -             & One-way efficiency ($\sqrt{0.9}$)                                                                 \\
        $\eta_{\text{out}}$         & 0.9487                     & -             & One-way efficiency ($\sqrt{0.9}$)                                                                 \\
        $SoC_{\min}$                & 2.00                       & MWh           & Minimum SoC health bound (10\% of $E_{\text{cap}}$)                                               \\
        $SoC_{\max}$                & 18.00                      & MWh           & Maximum SoC health bound (90\% of $E_{\text{cap}}$)                                               \\
        $SoC_{0}=0.5E_{\text{cap}}$ & 10.00                      & MWh           & Start at neutral SoC (50\% of $E_{\text{cap}}$)                                                   \\
        $C_{\text{deg}}$            & 25.0                       & €/MWh         & Conservative marginal throughput wear cost                                                        \\
        $\Delta t$                  & 1.0                        & h             & Modeling choice; hourly optimization granularity                                                  \\
        $P_{\text{bid,min}}$        & 1.0                        & MW            & German market specification                                                                       \\
        $\Delta P_{\text{bid}}$     & 1.0                        & MW            & German market specification                                                                       \\
        \midrule
        Terminal boundary choice    & Lower-bounded terminal SoC & -             & Fixed as $SoC_0 = 0.5E_{\text{cap}}$ and $SoC_T \ge 0.5E_{\text{cap}}$                            \\
        \bottomrule
    \end{tabularx}
\end{table}

The baseline parameterization uses $P_{\text{max}} = 10\,\text{MW}$ and $E_{\text{cap}} = 20\,\text{MWh}$, which corresponds to typical commercial grid-scale \gls{bess} in Germany suitable for frequency regulation markets (incl. \gls{afrr}). This configuration is validated by real-world projects such as Stadtwerke Mühlacker (10 MW / 20 MWh) \citep{Butler2026}, EnBW Marbach (15 MW / 22 MWh) \citep{EnBW2024}, and Qair Fuchsstadt (20 MW / 44 MWh) \citep{Qair2025}. These projects represent the common 1--2 hour duration class for utility-scale \gls{bess} participating in German electricity markets. \gls{soc} bounds (10--90\%) and initial conditions follow standard battery health management practices.

\subsubsection{Physical State and Feasibility Constraints}
Convention: $E_{\text{recv}, t}$ and $E_{\text{sent}, t}$ are AC-side metered energies at the grid connection point; therefore, charging maps to cells via $\eta_{\text{in}}$ and discharging maps from cells via $\eta_{\text{out}}$.

\begin{equation}\label{eq:act-split}
    E_{\text{act}, t}=E_{\text{act}, t}^{+}-E_{\text{act}, t}^{-},
    \qquad
    E_{\text{act}, t}^{+}\ge 0,\;E_{\text{act}, t}^{-}\ge 0
\end{equation}
\begin{equation}\label{eq:act-split-calc}
    E_{\text{act}, t}^{+}=
    \begin{cases}
        E_{\text{act}, t}, & \text{if } E_{\text{act}, t}>0 \\
        0,                 & \text{otherwise}
    \end{cases}
    ,\qquad
    E_{\text{act}, t}^{-}=
    \begin{cases}
        -E_{\text{act}, t}, & \text{if } E_{\text{act}, t}<0 \\
        0,                  & \text{otherwise}
    \end{cases}
\end{equation}
\begin{equation}\label{eq:soc-flow-map}
    E_{\text{recv}, t}^{\text{SoC}}=E_{\text{recv}, t}+E_{\text{act}, t}^{-},
    \qquad
    E_{\text{sent}, t}^{\text{SoC}}=E_{\text{sent}, t}+E_{\text{act}, t}^{+}
\end{equation}
Formulas \ref{eq:act-split} and \ref{eq:act-split-calc} define how net signed activation is split into upward and downward components.
Formula \ref{eq:soc-flow-map} then adds these directional components to the physical \gls{soc} charge/discharge flows.
In the hourly net-mode model, activation is represented by net signed energy $E_{\text{act}, t}$; $E_{\text{act}, t}^{+}$ and $E_{\text{act}, t}^{-}$ are directional decomposition terms used for linear accounting.

\begin{equation}\label{eq:soc-dynamics}
    SoC_{t} =
    \begin{cases}
        SoC_{t-1} + (E_{\text{recv}, t}^{\text{SoC}} \cdot \eta_{\text{in}}),                 & \text{if charging } (E_{\text{recv}, t}^{\text{SoC}} > 0)    \\
        SoC_{t-1} - \left( \frac{E_{\text{sent}, t}^{\text{SoC}}}{\eta_{\text{out}}} \right), & \text{if discharging } (E_{\text{sent}, t}^{\text{SoC}} > 0) \\
        SoC_{t-1},                                                                            & \text{otherwise}
    \end{cases}
\end{equation}
Charging uses $E_{\text{recv}, t}^{\text{SoC}}\cdot\eta_{\text{in}}$ in Formula \ref{eq:soc-dynamics} because only the efficiency-adjusted fraction of grid energy is stored in the cells.
Discharging uses $E_{\text{sent}, t}^{\text{SoC}}/\eta_{\text{out}}$ in Formula \ref{eq:soc-dynamics} because exporting a given grid energy volume requires higher internal battery energy due to conversion losses.

\begin{equation}\label{eq:soc-bounds}
    SoC_{\min} \leq SoC_{t} \leq SoC_{\max}
\end{equation}
In Formula \ref{eq:soc-bounds}, this limit keeps the battery in a safe operating range and avoids harmful overcharging or deep discharging.

\begin{equation}\label{eq:soc-cyclic}
    SoC_{0} = 0.5 \cdot E_{\text{cap}},
    \qquad
    SoC_{T} \ge 0.5 \cdot E_{\text{cap}}
\end{equation}
In Formula \ref{eq:soc-cyclic}, the simulation starts at 50\% \gls{soc} and must not end below 50\%, so the optimizer cannot exploit horizon-end discharge while still allowing feasible reserve and dispatch decisions.

\subsubsection{Operational Power and Market-Configuration Constraints}
The implemented dispatch optimizer uses a continuous linear formulation. Charging and discharging powers are non-negative continuous variables and are bounded independently by the converter rating $P_{\text{max}}$. The formulation therefore does not introduce a binary charge/discharge mode switch.
\begin{equation}\label{eq:charging_limit}
    0 \le P_{\text{charge},t} \le P_{\max}
\end{equation}
\begin{equation}\label{eq:discharging_limit}
    0 \le P_{\text{discharge},t} \le P_{\max}
\end{equation}

Reserve commitments are coupled to the same converter limit. Let $R_t^{+}$ denote upward reserve and $R_t^{-}$ downward reserve. Upward reserve requires discharge headroom, while downward reserve requires charge headroom:
\begin{equation}\label{eq:reserve-power-coupling}
    P_{\text{charge},t} + R^{-}_{t} \le P_{\max}
\end{equation}
\begin{equation}\label{eq:reserve-power-coupling-pos}
    P_{\text{discharge},t} + R^{+}_{t} \le P_{\max}
\end{equation}
\begin{equation}\label{eq:reserve-total-limit}
    R^{+}_{t}+R^{-}_{t}\le R_{\max}.
\end{equation}
These constraints limit total converter use but do not impose strict charge/discharge exclusivity.

\begin{equation}\label{eq:bigm-charge-discharge}
    \begin{aligned}
        0 \leq E_{\text{recv}, t}^{\text{plan}} & \leq P_{\text{max}} \cdot \Delta t, \\
        0 \leq E_{\text{sent}, t}^{\text{plan}} & \leq P_{\text{max}} \cdot \Delta t
    \end{aligned}
    \quad \forall t.
\end{equation}
Formula \ref{eq:bigm-charge-discharge} gives the equivalent planned-energy bounds without a binary mode variable.

\begin{equation}\label{eq:market-mode}
    \texttt{multi}: \quad
    P_{\text{charge},t},\,P_{\text{discharge},t},\,R_t^{+},\,R_t^{-} \ge 0
\end{equation}
\begin{equation}\label{eq:market-mode-energy}
    \texttt{da\_only}: \quad
    R_t^{+}=R_t^{-}=0
\end{equation}
\begin{equation}\label{eq:market-mode-afrr}
    \texttt{afrr\_only}: \quad
    P_{\text{charge},t}=P_{\text{discharge},t}=0
\end{equation}
Formulas \ref{eq:market-mode}--\ref{eq:market-mode-afrr} describe run-configuration choices rather than hourly binary market-selection decisions. In \texttt{multi}, \gls{da} and \gls{afrr} variables are allowed simultaneously subject to converter and \gls{soc} constraints. In \texttt{da\_only}, reserve variables are clamped to zero. In \texttt{afrr\_only}, \gls{da} charge and discharge variables are clamped to zero.

\begin{equation}\label{eq:energy-feasibility-up}
    SoC_{t} - \frac{P_{\text{bid}, t}^{\uparrow} \cdot \Delta t}{\eta_{\text{out}}} \geq SoC_{\min}
\end{equation}
In Formula \ref{eq:energy-feasibility-up}, the battery must have enough energy to deliver the full upward reserve when called.

\begin{equation}\label{eq:energy-feasibility-down}
    SoC_{t} + P_{\text{bid}, t}^{\downarrow} \cdot \Delta t \cdot \eta_{\text{in}} \leq SoC_{\max}
\end{equation}
In Formula \ref{eq:energy-feasibility-down}, the battery must keep enough empty space to absorb the full downward reserve when called.
In the stochastic rolling-horizon implementation, these deterministic bounds are replaced by the soft-penalty chance-constraints defined in Eq. \ref{eq:soft-p90-pos} and \ref{eq:soft-p90-neg} to ensure feasibility during extreme activation events.

\begin{equation}\label{eq:activation-feasibility}
    0 \leq E_{\text{act}, t}^{+} \leq P_{\text{bid}, t}^{\uparrow}\cdot\Delta t,
    \qquad
    0 \leq E_{\text{act}, t}^{-} \leq P_{\text{bid}, t}^{\downarrow}\cdot\Delta t
\end{equation}
Formula \ref{eq:activation-feasibility} limits upward and downward activated energy separately by the accepted directional reserve bids over the interval length.

\begin{equation}\label{eq:bid-bigm}
    \begin{aligned}
        0 \leq P_{\text{bid}, t}^{\uparrow}   & \leq P_{\text{bid,max}},     \\
        0 \leq P_{\text{bid}, t}^{\downarrow} & \leq P_{\text{bid,max}}.
    \end{aligned}
\end{equation}
Formula \ref{eq:bid-bigm} bounds directional reserve bids as continuous variables.

\begin{equation}\label{eq:bid-granularity}
    0 \leq P_{\text{bid},t}^{d} \leq P_{\text{bid,max}} = P_{\text{max}},
    \qquad d\in\{\uparrow,\downarrow\}.
\end{equation}
Formula \ref{eq:bid-granularity} states the continuous bid-volume bound used in the optimization; any reporting or market-facing rounding is handled outside the LP decision variables.

\subsubsection{Day-Ahead \& Intraday Cashflow Model}
To avoid pricing inconsistencies, \gls{da} and \gls{id} rescue trades are settled with distinct price processes. Let
$E^{DA}_{\mathrm{recv},t},E^{DA}_{\mathrm{sent},t}$ denote \gls{da} charging/discharging grid energies, and
$E^{ID}_{\mathrm{recv},t},E^{ID}_{\mathrm{sent},t}$ denote \gls{id} rescue charging/discharging grid energies.

\begin{equation}\label{eq:pi-da}
    \pi^{DA}_t
    =
    \underbrace{E^{DA}_{\mathrm{sent},t}\,\lambda_{\text{mkt},t}
        -
        E^{DA}_{\mathrm{recv},t}\,\lambda_{\text{mkt},t}}_{\text{DA energy cashflow}}
\end{equation}

\begin{equation}\label{eq:id-prices}
    p^{ID}_{\mathrm{buy},t}=\min\!\left(3000,\lambda_{\text{mkt},t}+30\right),
    \qquad
    p^{ID}_{\mathrm{sell},t}=\max\!\left(-500,\lambda_{\text{mkt},t}-30\right)
\end{equation}

\begin{equation}\label{eq:pi-id}
    \pi^{ID}_t
    =
    \underbrace{E^{ID}_{\mathrm{sent},t}\,p^{ID}_{\mathrm{sell},t}
        -
        E^{ID}_{\mathrm{recv},t}\,p^{ID}_{\mathrm{buy},t}}_{\text{ID rescue energy cashflow}}
\end{equation}

\begin{equation}\label{eq:pi-daid-realized-separated}
    \pi^{\mathrm{DA/ID}}_t
    =
    \pi^{DA}_t+\pi^{ID}_t
    -
    \underbrace{C^{\mathrm{deg,DA/ID}}_t}_{\text{degradation on internal throughput}}
\end{equation}

with
\begin{equation}\label{eq:cdeg-daid}
    C^{\mathrm{deg,DA/ID}}_t
    =
    C_{\mathrm{deg}}
    \left(
    \eta_{\text{in}}\!\left(E^{DA}_{\mathrm{recv},t}+E^{ID}_{\mathrm{recv},t}\right)
    +
    \frac{E^{DA}_{\mathrm{sent},t}+E^{ID}_{\mathrm{sent},t}}{\eta_{\text{out}}}
    \right).
\end{equation}

\subsubsection{aFRR Ex-Post Settlement Model}


\subsubsection{Capacity Revenue ($\pi_{t}^{\text{cap}}$)}

\begin{equation}\label{eq:pi-cap}
    \pi_{t}^{\text{cap}} = P_{\text{bid}, t} \cdot C_{\text{bid}, t} \cdot \Delta t
\end{equation}
Formula \ref{eq:pi-cap} calculates the reserve-availability payment: reserved MW multiplied by the accepted capacity price and the length of the time interval.


\subsubsection{Activation Revenue ($\pi_{t}^{\text{act}}$)}

The model uses signed activation energy values ($E_{\text{act}, t}$). In the German \gls{afrr} settlement representation used here, $\lambda_{\text{clr}, t}$ is the signed clearing price paid by the \gls{tso} to the provider. Therefore, $E_{\text{act}, t}\cdot\lambda_{\text{clr}, t}$ captures both directions: upward activation ($E_{\text{act}, t}>0$) yields revenue, while downward activation ($E_{\text{act}, t}<0$) yields a cost unless $\lambda_{\text{clr}, t}<0$.

Cost separation is defined to avoid double counting: ex-post activation cashflow includes settlement revenue and non-energy costs only ($C_{\text{pos}, t}$, $C_{\text{neg}, t}$), while energy replenishment effects are monetized later through realized \gls{da}/\gls{id} trades.

\begin{equation}\label{eq:c-pos}
    C_{\text{pos}, t} =
    \underbrace{\left( \frac{E_{\text{act}, t}}{\eta_{\text{out}}} \cdot C_{\text{deg}} \right)}_{\text{Degradation (discharge)}}
    \quad \text{for upward activation } (E_{\text{act}, t} > 0)
\end{equation}

\begin{equation}\label{eq:c-neg}
    C_{\text{neg}, t} =
    \underbrace{\left( \lvert E_{\text{act}, t}\rvert \cdot \eta_{\text{in}} \cdot C_{\text{deg}} \right)}_{\text{Degradation (charge)}}
    \quad \text{for downward activation } (E_{\text{act}, t} < 0)
\end{equation}
Formulas \ref{eq:c-pos} and \ref{eq:c-neg} define non-energy activation costs as direction-specific degradation wear.

\begin{equation}\label{eq:pi-act}
    \pi_{t}^{\text{act}} =
    \begin{cases}
        (E_{\text{act}, t} \cdot \lambda_{\text{clr}, t}) - C_{\text{pos}, t}, & \text{if upward activation } (E_{\text{act}, t} > 0)   \\
        (E_{\text{act}, t} \cdot \lambda_{\text{clr}, t}) - C_{\text{neg}, t}, & \text{if downward activation } (E_{\text{act}, t} < 0) \\
        0,                                                                     & \text{otherwise}
    \end{cases}
\end{equation}
Formula \ref{eq:pi-act} calculates activation profit from the signed activation payment and subtracts activation fees and degradation costs. Any later energy refill is reflected in the realized \gls{da}/\gls{id} cashflows.

\subsubsection{Total aFRR Ex-Post Profit ($\pi_{t}^{\text{aFRR}}$)}


\begin{equation}\label{eq:pi-afrr}
    \pi_{t}^{\text{aFRR}} = \pi_{t}^{\text{cap}} + \pi_{t}^{\text{act}}
\end{equation}
Formula \ref{eq:pi-afrr} aggregates capacity and activation components into total realized \gls{afrr} profit.

\subsubsection{aFRR Ex-Ante Expected Profit Model}

\subsubsection{Forecast Mapping Definitions}
\begin{equation}\label{eq:map-c-bid}
    \hat{C}_{\text{bid}, t} = g_{\text{bid}}(\mathbf{X}_{t})
\end{equation}
Formula \ref{eq:map-c-bid} defines predicted capacity price as an \gls{ml} model output conditioned on feature vector $\mathbf{X}_{t}$.

\begin{equation}\label{eq:map-lambda-clr}
    \hat{\lambda}_{\text{clr}, t} = g_{\text{clr}}(\mathbf{X}_{t})
\end{equation}
Formula \ref{eq:map-lambda-clr} defines the forecasted \gls{afrr} activation clearing price from the same feature set.

\begin{equation}\label{eq:map-p-act}
    \hat{p}_{\text{acc}, j, t}^{d} = g_{d}\big(\mathbf{X}_{t}, m_{j}^{d}\big)
\end{equation}
Formula \ref{eq:map-p-act} means the classifier maps contemporaneous market features and the interval-specific spread midpoint to the probability that a bid in interval $j$ and direction $d$ is accepted.



\subsubsection{Intraday Rebalancing and Reactive Rescue Mechanism}

Because high-frequency continuous intraday order-book/tick data were not available, a synthetic and deterministic \gls{id} rescue mechanism was introduced to preserve physical feasibility and to reduce extreme \gls{afrr} non-delivery events. The mechanism is explicitly \emph{reactive} and not part of the optimization model. In particular, the rolling-horizon optimizer remains blind to \gls{id} execution and is calibrated with severe \gls{soc}-violation penalties (denoted by $\lambda$) in its chance-constraint penalty terms. This deliberate over-conservatism ensures methodological identification of forecasting value (including the quantile sweep) and prevents the optimizer from exploiting a synthetic \gls{id} layer for ex-ante speculative arbitrage (i.e., no ``free lunch'' channel is provided to the planning model).

Causality is enforced with a one-hour execution lag. At hour $t$, the requirement for rebalancing is assessed from the post-clearing state and expected reserve obligations for hour $t+1$; the corresponding \gls{id} trade is then executed physically in hour $t+1$. Hence, no same-hour retroactive correction is permitted.
\begin{equation}\label{eq:id-one-hour-lag}
    E^{ID}_{\mathrm{exec},t+1} = f\left(SoC_t^{\mathrm{post}},\,E_{\text{act},t},\,P_{\text{bid},t+1}^{\uparrow},\,P_{\text{bid},t+1}^{\downarrow}\right)
\end{equation}
where $SoC_t^{\mathrm{post}}$ denotes the post-clearing/post-settlement state at the end of hour $t$, and $f(\cdot)$ maps this state and the next-hour reserve commitments to a feasible rescue energy decision subject to inverter and \gls{soc} limits. To emulate urgency and liquidity costs, \gls{id} rescue prices are linked to the day-ahead reference price $\lambda_{\text{mkt},t}$ via an asymmetric premium and capped by technical exchange bounds:
\begin{align}
    p^{ID}_{\mathrm{buy},t}  & = \min\!\left(3000,\; \lambda_{\text{mkt},t}+30\right), \\
    p^{ID}_{\mathrm{sell},t} & = \max\!\left(-500,\; \lambda_{\text{mkt},t}-30\right).
\end{align}
These rules impose a structural slippage wedge and avoid mathematically invalid prices in stress periods.

Physical realizability is enforced by inverter headroom after \gls{da} dispatch. Let $P_{\text{max}}$ denote converter power, and $P^{DA}_{\mathrm{ch},t},P^{DA}_{\mathrm{dis},t}$ the cleared \gls{da} charging/discharging powers. Then \gls{id} rescue power is bounded by
\begin{equation}
    P^{ID}_{\mathrm{ch},t} \le P_{\text{max}}-P^{DA}_{\mathrm{ch},t}, \qquad
    P^{ID}_{\mathrm{dis},t} \le P_{\text{max}}-P^{DA}_{\mathrm{dis},t}.
\end{equation}
Thus, energy is not ``teleported'' across markets; all rescue flows must pass through the same hardware bottleneck.

Degradation costs for \gls{id} rescue flows are aggregated with \gls{da} flows based on internal chemical throughput, as formalized in Equation~\ref{eq:cdeg-daid}.


\subsubsection{Machine Learning Target: Probabilistic Acceptance}
Instead of a continuous regression target for activated energy, the learning problem is framed as a probabilistic classification task. For each hour $t$ and direction $d\in\{\uparrow,\downarrow\}$, the model outputs the discrete probability vector $\big\{\hat{p}_{\text{acc}, j, t}^{d}\big\}_{j=1}^{J}$ over the set of spread intervals $j\in\{1,\dots,J\}$. This target aligns model predictions with the bid-acceptance mechanism of the \gls{afrr} merit order, allowing direct integration of data-driven acceptance likelihoods into the stochastic profit maximization.

\subsubsection{Interval construction (Data Preprocessing)}
The goal is to avoid guessing one exact bid price in a volatile market. Instead, the pricing landscape is sliced into discrete “zones” that are easier to reason about and optimize over. Each zone is defined by a bid spread, which represents how much higher or lower the bid is relative to the \gls{da} price. Using historical data, the continuous bid spread (historical bid minus historical \gls{da} price) is partitioned into $J$ uneven intervals $j\in\{1,\dots,J\}$. Quantile binning creates narrow intervals where the market is crowded (near the \gls{da} price) and wider intervals in the sparse tails, capturing the highly non-linear Residual Demand Curve \cite{aneiros2013functional} while preserving tractability for linear optimization \cite{abate2025bidding}. This turns a difficult continuous pricing problem into a manageable set of spread “zones” \cite{nolzen2022market}.

\subsubsection{Probabilistic Prediction (Machine Learning)}
Here the model behaves like a risk assessor rather than a single-point forecaster. For every hour $t$ and each spread interval $j$, it estimates the chance that the grid operator accepts the bid in a given direction $d\in\{\uparrow,\downarrow\}$. Concretely, the classifier outputs $\hat{p}_{\text{acc}, j, t}^{d}$, the predicted acceptance probability for zone $j$. This probability tells us how likely a bid at that price zone is to clear, given current market conditions \cite{koechlin2025strategic}. With a full probability vector across all zones, the optimizer can balance risk and reward rather than relying on a single brittle point forecast.

\subsubsection{Expected Activation Profit Formulation}
The battery’s energy is scarce: selling (upward activation) or buying (downward activation) in \gls{afrr} means forgoing the same transaction in the \gls{da} market. Profit must therefore be measured as an \emph{opportunity margin} relative to the \gls{da} price.

\textbf{Directional intuition.}
\begin{itemize}
    \item Upward activation ($s^{\uparrow}=1$): the battery discharges and sells energy; it earns a premium over what it could have received in the \gls{da} market.
    \item Downward activation ($s^{\downarrow}=-1$): the battery charges by buying energy; it gains a saving compared with what it would have paid in the \gls{da} market.
    \item $C_{\text{marg}, t}^{d}$ captures marginal degradation and efficiency losses for that direction.
\end{itemize}

\textbf{Pricing step.}
Let $m_{j}^{d}$ be the midpoint (€/MWh) of spread interval $j$ in direction $d$. The absolute bid price follows directly:
\[
    \hat{\lambda}_{\text{bid}, j, t}^{d} = \hat{\lambda}_{\text{mkt}, t} + m_{j}^{d}.
\]

\textbf{Opportunity margin and expected profit.}
Using the directional sign $s^{d}$, the opportunity margin simplifies to $\big(s^{d}\cdot m_{j}^{d} - C_{\text{marg}, t}^{d}\big)$. The expected activation profit for interval $j$ is then
\begin{equation}\label{eq:pi-act-prob}
    \hat{\pi}_{t}^{\text{act},d}(j) = \hat{p}_{\text{acc}, j, t}^{d} \cdot \Big( s^{d} \cdot m_{j}^{d} - C_{\text{marg}, t}^{d} \Big) \cdot P_{\text{bid}, t}^{d} \cdot \Delta t
\end{equation}
with $P_{\text{bid}, t}^{d}$ the bid volume (MW).
% EXPLANATION FOR WHY THIS IS A SMART CALCULATION 
% \paragraph{Derivation of the Directional Profit Margin}
% The formulation introduces a directional sign indicator $s^{d} \in \{-1, 1\}$ to consolidate the expected profit calculations for both upward and downward activations into a single equation. 

% First, recall the definition of the spread midpoint $m_{j}^{d}$, which represents the price spread of the absolute bid price $\hat{\lambda}_{\text{bid}, j, t}^{d}$ relative to the forecasted Day-Ahead market price $\hat{\lambda}_{\text{mkt}, t}$:
% \begin{equation}
%     m_{j}^{d} = \hat{\lambda}_{\text{bid}, j, t}^{d} - \hat{\lambda}_{\text{mkt}, t}
% \end{equation}

% Because the aFRR market involves two opposite physical and financial actions, the opportunity margin must be evaluated directionally:

% \begin{itemize}
%     \item \textbf{Case 1: Upward Activation ($d = \uparrow$).} 
%     The battery discharges, selling energy to the TSO instead of the Day-Ahead market. The profit margin is the premium earned over the Day-Ahead price:
%     \begin{equation*}
%         \text{Margin}^{\uparrow} = \hat{\lambda}_{\text{bid}, j, t}^{\uparrow} - \hat{\lambda}_{\text{mkt}, t} = m_{j}^{\uparrow}
%     \end{equation*}
%     By defining the upward sign indicator as $s^{\uparrow} = 1$, the term $s^{\uparrow} \cdot m_{j}^{\uparrow}$ evaluates exactly to $m_{j}^{\uparrow}$.

%     \item \textbf{Case 2: Downward Activation ($d = \downarrow$).} 
%     The battery charges, buying energy from the TSO instead of the Day-Ahead market. The profit margin in this case is the savings achieved compared to buying on the Day-Ahead market:
%     \begin{equation*}
%         \text{Margin}^{\downarrow} = \hat{\lambda}_{\text{mkt}, t} - \hat{\lambda}_{\text{bid}, j, t}^{\downarrow} = -\Big( \hat{\lambda}_{\text{bid}, j, t}^{\downarrow} - \hat{\lambda}_{\text{mkt}, t} \Big) = -m_{j}^{\downarrow}
%     \end{equation*}
%     By defining the downward sign indicator as $s^{\downarrow} = -1$, the term $s^{\downarrow} \cdot m_{j}^{\downarrow}$ evaluates exactly to $-m_{j}^{\downarrow}$.
% \end{itemize}

% Using this directional sign $s^{d}$, the gross opportunity margin seamlessly adapts to both directions as $s^{d} \cdot m_{j}^{d}$. Finally, subtracting the physical directional marginal cost $C_{\text{marg}, t}^{d}$ (accounting for degradation and efficiency losses) yields the consolidated net profit margin:
% \begin{equation*}
%     \Big( s^{d} \cdot m_{j}^{d} - C_{\text{marg}, t}^{d} \Big)
% \end{equation*}
% This algebraic substitution removes the need to define piecewise objective functions for charging and discharging.


\textbf{Choosing the best zone.}
Because the upward and downward books clear independently, the optimizer selects the best spread interval in each direction:
\begin{equation}\label{eq:pi-act-max}
    j^{*,d}(t) = \arg\max_{j\in\{1,\dots,J\}} \hat{\pi}_{t}^{\text{act},d}(j) \quad \forall d \in \{\uparrow, \downarrow\}
\end{equation}
yielding the interval that maximizes expected activation profit for that hour and direction.

\subsubsection{Directional Marginal Costs ($C_{\text{marg}, t}^{\uparrow}, C_{\text{marg}, t}^{\downarrow}$)}

The model uses direction-specific internal thresholds. Upward activation (discharging) and downward activation (charging) have different marginal costs because efficiency enters with opposite direction.


\begin{equation}\label{eq:c-marg-up}
    C_{\text{marg}, t}^{\uparrow} = \left( \frac{1}{\eta_{\text{out}}} \cdot C_{\text{deg}} \right)
\end{equation}

\begin{equation}\label{eq:c-marg-down}
    C_{\text{marg}, t}^{\downarrow} = \left( \eta_{\text{in}} \cdot C_{\text{deg}} \right)
\end{equation}
These deterministic marginal cost thresholds define the minimum activation price levels at which activation is economically worthwhile for the battery (see Formulas \ref{eq:c-marg-up} and \ref{eq:c-marg-down}).


\subsubsection{Total Expected aFRR Profit ($\hat{\pi}_{t}^{\text{aFRR}}$)}

\begin{equation}\label{eq:e-pi-afrr}
    \hat{\pi}_{t}^{\text{aFRR}} =
    \underbrace{P_{\text{bid}, t} \cdot \hat{C}_{\text{bid}, t} \cdot \Delta t}_{\text{Expected Capacity Revenue}}
    +
    \underbrace{\hat{\pi}_{t}^{\text{act}}}_{\text{Expected Activation Revenue}}
\end{equation}
Formula \ref{eq:e-pi-afrr} combines expected capacity revenue and the optimized expected activation profit into the total ex-ante \gls{afrr} value.

\subsubsection{Master Optimization Objective}

\begin{equation}\label{eq:objective}
    \max \ \hat{\Pi}_{\text{tot}} = \sum_{t=1}^{T} \left(\hat{\pi}_{t}^{\text{DA/ID}} + \hat{\pi}_{t}^{\text{aFRR}}\right)
\end{equation}
This keeps both terms in the optimization objective in Formula \ref{eq:objective} in ex-ante form and avoids mixing realized and expected quantities.

\subsubsection{Ex-Post Total Realized Profit}
\begin{equation}\label{eq:objective-realized}
    \Pi^{\mathrm{real}} = \sum_{t=1}^{T}\left(\pi_{t}^{\text{DA/ID}} + \pi_{t}^{\text{aFRR}} - C_{pen,t} - C_{tx,t} - C^{\mathrm{deg}}_t\right) + V_{\mathrm{term}}
\end{equation}
In Formula \ref{eq:objective-realized} the total ex-post/realized profit is calculated.

\subsubsection{Model Assumptions and Scope}
The proposed formulation is a simplified representation of \gls{bess} trading behavior under operational and market constraints.

\textbf{Scope and market convention (DE/LU).}
The \gls{afrr} activation term is implemented with signed quantities, i.e., $E_{\text{act}, t}\cdot\lambda_{\text{clr}, t}$, according to the \gls{delu} market-rule representation used in this thesis.
This is a scope and settlement-mapping choice for the selected market, not a free modeling assumption.

\textbf{Assumption: electrochemical degradation proxy.}
Degradation is modeled as a fixed cost per internal MWh that passes through the battery cells.
This keeps the model simple and interpretable, but it does not capture detailed aging mechanisms such as temperature effects, calendar aging, or nonlinear state-of-health changes.
Therefore, $C_{\text{deg}}$ should be read as an average wear cost for the selected battery technology and operating pattern.

\textbf{Assumption: auxiliary and transaction costs.}
Auxiliary consumption is neglected because, for the parameterization used here, it is not large enough to materially change the trading decisions or economic results.
Transaction costs are retained as a linear accounting term $C_{tx,t}$ and can be set to zero in baseline runs or varied in sensitivity analyses.
Accordingly, the model focuses on the dominant trading cashflows, penalty exposure, and degradation costs.

\subsubsection{Methodological and Reporting Choices}
\textbf{Temporal granularity and dispatch structure.}
The model uses a fixed hourly resolution ($\Delta t = 1$ h).
Although the \gls{da} market moved to quarter-hourly intervals on October 1, 2025, the dataset used here starts in 2022 and is predominantly hourly; therefore, a uniform hourly resolution is used for consistency across the full sample.
This choice abstracts from sub-hourly activation dynamics, which can be relevant in balancing markets.
The dispatch problem is solved as a continuous LP relaxation. This improves computational robustness, but it also means that charge/discharge exclusivity is not enforced by a binary mode switch.

\textbf{Modeling of the intraday market interface.}

\gls{id} rebalancing is modeled via the continuous reactive rescue mechanism formally defined in Subsection~\ref{subsec:dispatch-optimization} and in the dedicated subsection on intraday rebalancing above. In methodological terms, this interface acts as an ex-post feasibility safeguard under strict causality (with execution lag), not as an ex-ante optimization channel. Consequently, the LP planner remains blind to intraday rescue in order to prevent synthetic arbitrage and preserve identification of forecasting value in the \gls{da}/\gls{afrr} decision layer.

\textbf{Ex-ante uncertainty modeling.}
Expected activation profits are based on interval-level bid-acceptance probabilities $\hat{p}_{\text{acc}, j, t}^{d}$ and their associated spread midpoints $m_{j}^{d}$.
This discretized, probabilistic formulation aligns model outputs with the \gls{afrr} merit-order acceptance process and avoids imposing deterministic full-hour activation assumptions.
The predictive validity of the ex-ante objective depends on the calibration quality of the probabilistic classifier and on stable feature-to-market relationships across the spread bins.

\textbf{Accounting convention and interpretation.}
To avoid accounting overlap, the ex-post formulation includes non-energy activation costs only ($C_{\text{pos}, t}$, $C_{\text{neg}, t}$), while energy replenishment is valued through subsequent realized \gls{da}/\gls{id} trades.
Results should be interpreted as conditional on these assumptions and modeling choices; robustness checks should include sensitivity analyses for $C_{\text{deg}}$, efficiency parameters, and the calibration of the interval-level acceptance probabilities.


\subsubsection{Ex-Post Capacity Non-Delivery Penalties}
Capacity non-delivery is evaluated on post-rescue feasibility and is jointly constrained by (i) available \gls{soc} energy headroom/footroom and (ii) residual inverter capacity after \gls{da} and \gls{id} execution.

\begin{equation}\label{eq:rmax-pos}
    R^{+}_{\max,t}
    =
    \min\!\left(
    \frac{(SoC_t-SoC_{\min})\eta_{\text{out}}}{\Delta t},
    \;
    P_{\text{max}}-P^{DA}_{\mathrm{sent},t}-P^{ID}_{\mathrm{sent},t}
    \right)
\end{equation}

\begin{equation}\label{eq:rmax-neg}
    R^{-}_{\max,t}
    =
    \min\!\left(
    \frac{SoC_{\max}-SoC_t}{\eta_{\text{in}}\Delta t},
    \;
    P_{\text{max}}-P^{DA}_{\mathrm{recv},t}-P^{ID}_{\mathrm{recv},t}
    \right)
\end{equation}

Missed reserve capacity is
\begin{equation}\label{eq:missed-cap}
    M^{+}_{t}=\max\!\left(0, P_{\text{bid},t}^{\uparrow}-R^{+}_{\max,t}\right),
    \qquad
    M^{-}_{t}=\max\!\left(0, P_{\text{bid},t}^{\downarrow}-R^{-}_{\max,t}\right).
\end{equation}

Activation non-delivery is measured by comparing requested and physically delivered activation energy:
\begin{equation}\label{eq:missed-activation}
    E^{act,+}_{miss,t}=\max\!\left(0,\,E^{act,+}_{req,t}-E^{act,+}_{del,t}\right),
    \qquad
    E^{act,-}_{miss,t}=\max\!\left(0,\,E^{act,-}_{req,t}-E^{act,-}_{del,t}\right).
\end{equation}
Here, $E^{act,+}_{req,t}$ and $E^{act,-}_{req,t}$ denote requested upward and downward activation energy, while $E^{act,+}_{del,t}$ and $E^{act,-}_{del,t}$ denote delivered activation energy after physical and \gls{id} rescue constraints.

At settlement level, the official incentive logic for non-delivery can be represented as \citep{Regelleistung.net2026}:
\begin{equation}\label{eq:cpen-official}
    C^{\mathrm{official}}_{\mathrm{pen},t}
    =
    M_t \cdot \max\!\left(\overline{C}_{\mathrm{cap, slice}},\, \lambda_{\mathrm{ID, AEP}} + \delta\right),
\end{equation}
where $M_t$ denotes the unavailable reserve capacity, $\overline{C}_{\mathrm{cap, slice}}$ is the average capacity price of the relevant product slice, $\lambda_{\mathrm{ID, AEP}}$ is the \gls{id} AEP reference, and $\delta$ is the regulatory surcharge.

The simulation decomposes total non-delivery penalties into activation and capacity components:
\begin{equation}\label{eq:cpen-total}
    C^{pen}_t=C^{pen,act}_t+C^{pen,cap}_t.
\end{equation}
Activation penalties are calculated from missed activation energy and a direction-specific penalty basis:
\begin{equation}\label{eq:cpen-act}
    C^{pen,act}_t
    =
    E^{act,+}_{miss,t}\,B^{act,+}_t
    +
    E^{act,-}_{miss,t}\,B^{act,-}_t,
\end{equation}
\begin{equation}\label{eq:cpen-act-basis}
    B^{act,+}_t=\max\!\left(|\lambda^{marg,+}_t|,\,|\lambda^{IDAEP,+}_t+\delta^{act}|\right),
    \qquad
    B^{act,-}_t=\max\!\left(|\lambda^{marg,-}_t|,\,|\lambda^{IDAEP,-}_t+\delta^{act}|\right).
\end{equation}
Capacity penalties are calculated from unavailable reserve capacity and a direction-specific capacity penalty basis:
\begin{equation}\label{eq:cpen-cap}
    C^{pen,cap}_t
    =
    M^{+}_{t}\,B^{cap,+}_t\,\Delta t
    +
    M^{-}_{t}\,B^{cap,-}_t\,\Delta t,
\end{equation}
\begin{equation}\label{eq:cpen-cap-basis}
    B^{cap,+}_t=\max\!\left(\overline{C}_{\mathrm{cap, product}},\,\frac{|\lambda^{IDAEP,+}_t|}{1000}+\delta^{cap}\right),
    \qquad
    B^{cap,-}_t=\max\!\left(\overline{C}_{\mathrm{cap, product}},\,\frac{|\lambda^{IDAEP,-}_t|}{1000}+\delta^{cap}\right).
\end{equation}
The terms $B^{act,\pm}_t$ are expressed in \texteuro/MWh, while $B^{cap,\pm}_t$ are expressed in \texteuro/MW/h. If explicit marginal energy prices or \gls{id} AEP references are unavailable, the implementation uses configured fallback values to keep the backtest deterministic and reproducible.

\subsubsection{Financial Settlement and Liquidity}
Transaction costs are modeled proportional to the total grid-side energy throughput:
\begin{equation}\label{eq:ctx}
    C_{tx,t} = c_{tx} \cdot \left( E^{DA}_{\mathrm{recv},t} + E^{DA}_{\mathrm{sent},t} + E^{ID}_{\mathrm{recv},t} + E^{ID}_{\mathrm{sent},t} + |E_{\text{act},t}| \right)
\end{equation}
where $c_{tx}$ is the specific transaction fee in \texteuro/MWh.

The net liquid cashflow isolates actual cash movements, strictly excluding non-cash depreciation such as battery degradation ($C^{\mathrm{deg}}_t$):
\begin{equation}
    CF_{net,t} = R_{DA,t} - C_{DA,t} + R_{cap,t} + R_{act,t} - C_{tx,t} - C_{pen,t}
\end{equation}

The maximum working capital requirement represents the peak cumulative drawdown:
\begin{equation}
    \text{max\_cap\_req} = \max\left(0,\,-\min_t \left(Cash_0 + \sum_{k=1}^{t} CF_{net,k} \right)\right)
\end{equation}

Assuming $\text{max\_cap\_req} > 0$, the return on capital is defined as:
\begin{equation}
    ROI_{cap} = \frac{\Pi^{\mathrm{real}}}{\text{max\_cap\_req}}
\end{equation}

\subsubsection{Terminal Valuation and Operational KPIs}
To account for unsold energy at the end of the simulation horizon ($T$), the remaining inventory strictly above the minimal state of charge ($SoC_{\min}$) is marked-to-market at the final \gls{da}/\gls{id} reference price, discounting inverter losses ($\eta_{\text{out}}$):
\begin{equation}
    V_{\mathrm{term}} = \max(0, (SoC_T - SoC_{\min})) \cdot \eta_{\text{out}} \cdot \lambda_{\text{mkt},T}
\end{equation}
The total realized profit is thus $\Pi^{\mathrm{real}} = \sum_t \pi^{\mathrm{real}}_t + V_{\mathrm{term}}$.


\subsection{Techno-Economic Assessment}
\label{subsec:techno-economic-assessment}
Techno-economic evaluation is performed on realized settlement outcomes and decomposed into value, risk, and operational stress indicators.

\paragraph{Value of Forecasting (VoF).}
For model $m$ and benchmark $b$, the incremental value of forecast quality is defined as:
\begin{equation}\label{eq:vof}
    VoF_{m\to b} = \Pi^{\mathrm{real}}_{m} - \Pi^{\mathrm{real}}_{b}.
\end{equation}
Relevant comparisons include probabilistic models against naive baselines and against deterministic alternatives. Positive \gls{vof} indicates economically beneficial forecast information under identical physical constraints.

\paragraph{Capture and benchmark ratios.}
To normalize across periods with different market scales, profit capture is evaluated relative to benchmark and oracle envelopes:
\begin{equation}
    CR^{\mathrm{bench}}_{m} = \frac{\Pi^{\mathrm{real}}_{m}}{\Pi^{\mathrm{real}}_{\mathrm{bench}}},
    \qquad
    CR^{\mathrm{oracle}}_{m} = \frac{\Pi^{\mathrm{real}}_{m}}{\Pi^{\mathrm{real}}_{\mathrm{oracle}}}.
\end{equation}
$CR^{\mathrm{oracle}}_{m}\in(-\infty,1]$ measures proximity to the informational upper bound.

\paragraph{Profit-component decomposition.}
Total realized value is reported as:
\begin{equation}
    \Pi^{\mathrm{real}} =
    \underbrace{\Pi^{DA/ID}}_{\text{energy trading}}
    +
    \underbrace{\Pi^{cap}}_{\text{aFRR capacity}}
    +
    \underbrace{\Pi^{act}}_{\text{aFRR activation}}
    -
    \underbrace{C^{\mathrm{tx}}}_{\text{transaction}}
    -
    \underbrace{C^{\mathrm{pen}}}_{\text{non-delivery}}
    -
    \underbrace{C^{\mathrm{deg}}}_{\text{degradation}}
    +
    \underbrace{V_{\mathrm{term}}}_{\text{terminal inventory}}.
\end{equation}
This decomposition identifies whether gains stem from directional market timing, reserve monetization, or reduced penalty exposure.

\paragraph{Operational intensity and stress.}
Operational stress is reported via Equivalent Full Cycles (EFC), defined consistently with the discharge-based convention used in the settlement layer. Let total grid-side discharged energy be
\[
    E^{\mathrm{grid}}_{\mathrm{dis,tot}}
    =
    \sum_t \left(E^{DA}_{\mathrm{sent},t}+E^{ID}_{\mathrm{sent},t}\right)
    +
    \sum_t E^{+}_{\mathrm{act},t},
\]
then internal discharged throughput equals
\[
    E^{\mathrm{int}}_{\mathrm{dis,tot}}
    =
    \frac{E^{\mathrm{grid}}_{\mathrm{dis,tot}}}{\eta_{\text{out}}}.
\]
Hence,
\begin{equation}\label{eq:efc-corrected}
    \mathrm{EFC}
    =
    \frac{E^{\mathrm{int}}_{\mathrm{dis,tot}}}{E_{\text{cap}}}
    =
    \frac{\left(\sum_t \left(E^{DA}_{\mathrm{sent},t}+E^{ID}_{\mathrm{sent},t}\right)+\sum_t E^{+}_{\mathrm{act},t}\right)/\eta_{\text{out}}}{E_{\text{cap}}}.
\end{equation}
This avoids double counting from bidirectional throughput and keeps EFC fully aligned with the discharge-only cycle definition.

\paragraph{Feasibility diagnostics.}
To detect hidden infeasibility pressure in soft-constrained optimization, hourly slack utilizations are monitored:
\begin{equation}
    S^{+}=\sum_t s_t^{+},\qquad S^{-}=\sum_t s_t^{-}.
\end{equation}
Large $S^{+}$/$S^{-}$ indicate persistent quantile-based reserve-feasibility stress and are interpreted jointly with penalty payments and realized \gls{soc} trajectories.
